%% kapa: exam=1 round=31 year=2020 period=1 form=A subject=economics subject_ko=경제학원론 seq=23 number=63
%% source: Q-Net 감정평가사 기출문제 (cst00309) artlSeq=5207219 file=extracted/1008/5207219/A형/A형/2020년도 제31회 감정평가사 1차 시험_1교시_A형.pdf
%% conversion: kordoc 4.24.51 -> tools/gen1.py (auto); review=proofread-sonnet; flags=none
\begin{question}{63}
경제성장모형에서 甲국의 총생산함수가 $Q=AL^{0.75}K^{0.25}$ 일 때, 옳지 않은 것은? (단, $Q$ 는 생산량, $L$ 은 노동량, $K$ 는 자본량, 시장은 완전경쟁시장이다.)
\begin{choices}
  \item 자본탄력성은 0.25이다.
  \item 노동분배율은 자본분배율보다 크다.
  \item A는 총요소생산성이다.
  \item 노동량, 자본량 및 총요소생산성이 각각 10\% 씩 증가하면 생산량은 10\% 증가한다.
  \item 총생산함수는 규모에 대한 수익 불변이다.
\end{choices}
\end{question}
\answer{4}
